%=============================================================================
% Section 1: Introduction
%=============================================================================
\section{Introduction}

\subsection{The Limitation of Intelligence Alone}
Artificial intelligence systems have achieved remarkable progress in reasoning, language understanding, and multimodal perception. Large-scale models from providers such as OpenAI~\cite{openai2023}, Anthropic~\cite{anthropic2024}, Google, and DeepSeek demonstrate increasingly sophisticated cognitive capabilities. However, a critical limitation remains:

\begin{quote}
\textit{Intelligence alone does not produce real-world impact.}
\end{quote}

For an AI system to perform meaningful tasks, it must interact with external systems. Examples include:
\begin{itemize}[leftmargin=*]
\item Retrieving information from knowledge bases
\item Invoking APIs to perform actions
\item Controlling devices in the physical world
\item Executing automation workflows
\item Interacting with digital services
\end{itemize}

In other words, \textbf{AI must connect its internal reasoning with external action}. This connection is what we define as the AI Contact Layer.

\subsection{The Missing Conceptual Layer}
While execution layers and runtime systems address implementation concerns, there exists a deeper conceptual question: What is the nature of the interface between AI intelligence and the external world?

This question has both philosophical and practical dimensions:
\begin{itemize}[leftmargin=*]
\item \textbf{Philosophical:} How do we understand the boundary between thought and action in AI systems?
\item \textbf{Practical:} How do we architect systems that reliably bridge reasoning and execution?
\end{itemize}

The Contact Layer concept addresses both dimensions, providing:
\begin{enumerate}[label=\arabic*.]
\item A conceptual framework for understanding AI-world interaction
\item An architectural model for implementing this interaction
\item A unifying perspective that connects execution infrastructure with broader system design
\end{enumerate}

\subsection{Relationship to the Execution Layer Paper}
This paper builds upon the AI Execution Layer architecture described in our companion work~\cite{wang2026execution}. The relationship between these works can be understood through the lens of the AI Action Layer:

\textbf{Action Layer Decomposition:} The Action Layer comprises two subcomponents:
\begin{itemize}[leftmargin=*]
\item \textbf{Contact Layer:} Routing, policy evaluation, provider selection (the ``control plane'')
\item \textbf{Execution Runtime:} API invocation, protocol translation, error handling (the ``data plane'')
\end{itemize}

\textbf{Complementary Focus:} The Execution Layer paper describes the \textit{runtime component}---how capabilities are invoked, translated, and normalized. This paper describes the \textit{contact component}---how routing decisions are made, policies evaluated, and providers selected.

Together, they form the complete Action Layer: one handles the mechanics of execution, the other handles the logic of routing.

\begin{table}[h]
\centering
\caption{Comparison of Execution Layer and Contact Layer Papers}
\begin{tabular}{lll}
\toprule
\textbf{Aspect} & \textbf{Execution Layer Paper~\cite{wang2026execution}} & \textbf{Contact Layer Paper (This Work)} \\
\midrule
Focus & Runtime mechanics & Conceptual framework \\
Question & How to execute? & Why mediate? \\
Component & Execution Runtime & Contact Layer \\
Layer & Data plane & Control plane \\
Key Constraint & Minimality & Intelligence-Action separation \\
\bottomrule
\end{tabular}
\end{table}

\subsection{Contributions}
The primary contributions are as follows:
\begin{enumerate}[label=\arabic*.]
\item \textbf{Conceptual Framework:} We define the Contact Layer as the interface between AI intelligence and the external world, distinguishing it from both cognitive processes and execution mechanisms.

\item \textbf{Theoretical Foundation:} We establish the Intelligence-Action distinction as a fundamental principle of AI-native system design.

\item \textbf{Architectural Integration:} We demonstrate how the Contact Layer relates to execution runtimes and orchestration systems, clarifying terminology boundaries.

\item \textbf{Practical Implications:} We show how this conceptual framework informs the design of capability-mediated AI systems.
\end{enumerate}
